\Askhsh{37087}{4}
\begin{ylh}{1}
    \item 3.6. Κριτήρια ισότητας ορθογώνιων τριγώνων
    \item 4.2. Τέμνουσα δύο ευθειών - Ευκλείδειο αίτημα
    \item 5.3. Ορθογώνιο
    \item 5.6. Εφαρμογές στα τρίγωνα
    \item 5.8. Το ορθόκεντρο τριγώνου.
\end{ylh}
% ===================================================================
%  Αρχείο: 37087.tex (Fragment)
%  Αυτόματη μετατροπή μέσω doc2latex.py
% ===================================================================

ΘΕΜΑ 4

Δίνεται ισόπλευρο τρίγωνο $ΑΒ\Gamma$ και τα ύψη του $BK$ και $\Gamma\Lambda$, τα οποία τέμνονται στο $I$. Αν $M$ και $N$ είναι τα μέσα των $IB$ και $I\Gamma$ αντίστοιχα, να αποδείξετε ότι:

\begin{alist}
\item Το $AI$ προεκτεινόμενο διέρχεται από το μέσο της πλευράς $B\Gamma$. \monades{10}
\item Το τετράπλευρο $M\Lambda KN$ είναι ορθογώνιο παραλληλόγραμμο. \monades{15}

\begin{center}
\begin{tikzpicture}[scale=1.2]
\tkzDefPoint(0,0){B}
\tkzDefPoint(4,0){C}
\tkzDefEquilateral(B,C) \tkzGetPoint{A}
\tkzDefPointBy[projection=onto A--C](B) \tkzGetPoint{K}
\tkzDefPointBy[projection=onto A--B](C) \tkzGetPoint{L}
\tkzInterLL(B,K)(C,L) \tkzGetPoint{I}
\tkzDefMidPoint(I,B) \tkzGetPoint{M}
\tkzDefMidPoint(I,C) \tkzGetPoint{N}
\tkzDefMidPoint(B,C) \tkzGetPoint{midBC}

\draw[l3] (A)--(B)--(C)--cycle;
\draw[l3] (B)--(K);
\draw[l3] (C)--(L);
\draw[l3] (A)--(midBC);
\draw[l3] (M)--(L)--(K)--(N)--cycle;

\tkzDrawPoints(A,B,C,K,L,I,M,N)
\tkzLabelPoint[above](A){$A$}
\tkzLabelPoint[below left](B){$B$}
\tkzLabelPoint[below right](C){$\Gamma$}
\tkzLabelPoint[right](K){$K$}
\tkzLabelPoint[left](L){$\Lambda$}
\tkzLabelPoint[above right](I){$I$}
\tkzLabelPoint[below](M){$M$}
\tkzLabelPoint[below](N){$N$}

\tkzMarkRightAngle(B,K,C)
\tkzMarkRightAngle(C,L,B)
\end{tikzpicture}
\end{center}
\end{alist}
